
% FYI, the following packages are preloaded with the document class:
% longtable, xcolor, graphicx, booktabs, caption, csquotes, hyperref,
% calc, listings, datetime2, siunitx, geometry, enumitem

%%%%%%%%%%%%%%%%%%%%%%%%%%%%%%%%%%%%%%%%%%%%%%%%%%%%%%%%%% extra packages
\usepackage{amsmath}		% the main package in the AMS-LATEX distribution
\usepackage{amsfonts}		% extended set of fonts for use in mathematics
\usepackage{amssymb}		% adds new symbols to be used in math mode
\usepackage{mathrsfs}		% math fonts, e.g., Laplace
\usepackage{float}			% provides the H float modifier option
\usepackage{multirow}		% tables \multirow command
\usepackage{subcaption}		% enables subfigures
\usepackage{lscape}			% for landscape mode
\usepackage{verbatim}		% new verbatim environment, \begin{comment}...\end{comment}, \verbatiminput
%\usepackage{pgfgantt}		% for Gantt charts, see https://ctan.org/pkg/pgfgantt

%add extra packages if needed here

%%%%%%%%%%%%%%%%%%%%%%%%%%%%%%%%%%%%%%%%%%%%%%%%%%%%%%%%%% temp packages
\usepackage{lipsum}						% for fake text
%\usepackage[textsize=tiny]{todonotes}   % enable To-do notes, use the option "disable" to hide all notes, usage \todo{}

%\usepackage{draftwatermark}			% prints a watermark overlay, uncomment if needed
%\SetWatermarkText{**DRAFT**}
%\SetWatermarkScale{1}
%\SetWatermarkColor[gray]{0.8}

%%%%%%%%%%%%%%%%%%%%%%%%%%%%%%%%%%%%%%%%%%%%%%%%%%%%%%%%%% settings
\AtBeginDocument{					% Rendered PDF metadata:
\hypersetup{pdftitle=\ttitle} 		% Sets the PDF title to your dissertation title
\hypersetup{pdfauthor=\authorname} 	% Sets the PDF author to your name
}

%%%%%%%%%%%%%%%%%%%%%%%%%%%%%%%%%%%%%%%%%%%%%%%%%%%%%%%%%% user macros

%....









% Pacotes adicionais necessários às tabelas, equações e diagrama do conteúdo.
% Nenhuma redefinição da apresentação da classe original.
\usepackage{array,tabularx,tikz}
\usetikzlibrary{arrows.meta,positioning,calc}
\newcolumntype{Y}{>{\raggedright\arraybackslash}X}
\newcommand{\grau}[1]{#1^{\circ}}
\newcommand{\sat}{\operatorname{sat}}
\newcommand{\diag}{\operatorname{diag}}
\newcommand{\R}{\mathbf{R}}
\newcommand{\vet}[1]{\boldsymbol{#1}}
\newcommand{\fonte}{\par\smallskip\footnotesize Fonte: elaboração própria.}

% Marcador para partes por completar (aparece a vermelho no PDF)
\newcommand{\acompletar}[1]{\textcolor{red}{\textit{[A completar: #1]}}}
\usetikzlibrary{shapes.geometric,fit,backgrounds}
\usepackage[european]{circuitikz}
